\chapter{Về Bản Thân}
\markboth{Về Bản Thân}{About Yourself}

\section{Không Ai Hoàn Hảo}
\begin{truyen}{Không Ai Hoàn Hảo}{No One Is Perfect}
\chuhoa{K}{hông ai học giỏi tất cả, cư xử đúng tất cả, và mạnh ở mọi mặt.}
\textit{(No one is excellent at everything, always behaves perfectly, or is strong in every area.)}

Hiểu điều này sớm giúp học sinh bớt tự ép mình theo một hình mẫu không có thật.
\textit{(Understanding this early helps students stop forcing themselves toward an unreal ideal.)}

Con người vẫn có thể rất đáng quý dù còn thiếu sót.
\textit{(A person can still be deeply valuable while remaining imperfect.)}

Biết mình không hoàn hảo không phải để buông xuôi, mà để sống thật hơn.
\textit{(Knowing you are not perfect is not an excuse to give up, but a way to live more honestly.)}
\end{truyen}

\begin{baihoc}
Chấp nhận mình không hoàn hảo là bước đầu để trưởng thành bình tĩnh hơn.
\textit{(Accepting imperfection is a first step toward calmer growth.)}
\end{baihoc}

\begin{ghinhoanh}
\textbf{Short English to remember:}

Imperfection does not cancel worth.

\end{ghinhoanh}

\begin{tuhocnhanh}
1. Vì sao hiểu rằng không ai hoàn hảo lại quan trọng? \textit{Why is it important to know that no one is perfect?}

2. Điều này giúp học sinh bớt điều gì? \textit{What does it help students reduce?}

3. Em có đang đòi mình quá hoàn hảo không? \textit{Are you demanding too much perfection from yourself?}
\end{tuhocnhanh}
\ngancach

\section{Ai Cũng Có Điểm Mạnh}
\begin{truyen}{Ai Cũng Có Điểm Mạnh}{Everyone Has Strengths}
\chuhoa{M}{ỗi học sinh có một chỗ mạnh nào đó, dù chưa chắc nó hiện ra rất sớm hoặc rất rõ.}
\textit{(Every student has some strength, even if it does not appear early or clearly.)}

Có người mạnh về tính toán, có người mạnh về ngôn ngữ, có người mạnh về kiên trì hay cách đối xử với người khác.
\textit{(Some are strong in calculation, others in language, persistence, or human interaction.)}

Nhìn ra điểm mạnh của mình là một việc cần thiết.
\textit{(Recognizing your strengths is necessary.)}

Nó giúp em có chỗ đứng vững hơn trước rất nhiều so sánh ngoài kia.
\textit{(It gives you firmer ground against comparison.)}
\end{truyen}

\begin{baihoc}
Biết điểm mạnh của mình không để kiêu ngạo, mà để đi đúng hơn.
\textit{(Knowing your strengths is not for arrogance, but for clearer direction.)}
\end{baihoc}

\begin{ghinhoanh}
\textbf{Short English to remember:}

Strengths can guide your way forward.

\end{ghinhoanh}

\begin{tuhocnhanh}
1. Vì sao ai cũng cần nhìn ra điểm mạnh của mình? \textit{Why should everyone see their strengths?}

2. Điểm mạnh có thể nằm ở những đâu? \textit{Where can strengths appear?}

3. Em thấy điểm mạnh nào của mình rõ nhất? \textit{What strength of yours feels clearest?}
\end{tuhocnhanh}
\ngancach

\section{Ai Cũng Có Điểm Yếu}
\begin{truyen}{Ai Cũng Có Điểm Yếu}{Everyone Has Weaknesses}
\chuhoa{C}{ũng như điểm mạnh, điểm yếu là điều rất bình thường ở con người.}
\textit{(Just like strengths, weaknesses are a normal part of being human.)}

Vấn đề không phải là có điểm yếu hay không, mà là em có chịu nhìn nó cho thật không.
\textit{(The issue is not whether you have weaknesses, but whether you face them honestly.)}

Người biết điểm yếu của mình thường học cách đi cẩn thận hơn.
\textit{(Those who know their weaknesses often learn to move more wisely.)}

Điểm yếu không phải bản án; nó chỉ là phần cần được hiểu và rèn thêm.
\textit{(A weakness is not a sentence; it is simply something to understand and train.)}
\end{truyen}

\begin{baihoc}
Nhìn đúng điểm yếu giúp mình thực tế hơn, không phải kém giá trị hơn.
\textit{(Seeing your weaknesses clearly makes you more realistic, not less worthy.)}
\end{baihoc}

\begin{ghinhoanh}
\textbf{Short English to remember:}

Weaknesses invite growth, not despair.

\end{ghinhoanh}

\begin{tuhocnhanh}
1. Vì sao điểm yếu là điều bình thường? \textit{Why are weaknesses normal?}

2. Điều quan trọng không phải là gì? \textit{What is not the real issue?}

3. Em có điểm yếu nào cần nhìn thẳng hơn? \textit{What weakness do you need to face more honestly?}
\end{tuhocnhanh}
\ngancach

\section{Mỗi Người Học Khác Nhau}
\begin{truyen}{Mỗi Người Học Khác Nhau}{Everyone Learns Differently}
\chuhoa{K}{hông phải ai cũng hiểu bài theo cùng một cách và cùng một nhịp.}
\textit{(Not everyone understands in the same way or at the same pace.)}

Có người cần nghe lại, có người cần viết ra, có người cần làm thử nhiều lần.
\textit{(Some need to hear it again, some need to write it, others need repeated practice.)}

Hiểu điều này giúp học sinh bớt so sánh máy móc với người khác.
\textit{(Understanding this helps students compare themselves less mechanically.)}

Nó cũng giúp em tìm được cách học hợp hơn với chính mình.
\textit{(It also helps you find methods that fit you better.)}
\end{truyen}

\begin{baihoc}
Học khác nhau không có nghĩa là học kém hơn; chỉ là con đường hiểu bài khác nhau.
\textit{(Learning differently does not mean learning worse; it simply means a different path to understanding.)}
\end{baihoc}

\begin{ghinhoanh}
\textbf{Short English to remember:}

Different methods can still lead to understanding.

\end{ghinhoanh}

\begin{tuhocnhanh}
1. Vì sao mỗi người học khác nhau? \textit{Why does each person learn differently?}

2. Hiểu điều này giúp ích gì? \textit{How does understanding this help?}

3. Cách học nào hiện hợp với em nhất? \textit{What way of learning fits you best right now?}
\end{tuhocnhanh}
\ngancach

\section{Chậm Không Phải Là Dốt}
\begin{truyen}{Chậm Không Phải Là Dốt}{Being Slow Is Not Being Stupid}
\chuhoa{N}{hiều học sinh học chậm và vội kết luận rằng mình dốt.}
\textit{(Many students learn slowly and quickly conclude that they are stupid.)}

Nhưng chậm có thể đến từ nhiều lý do: hổng nền, chưa hợp cách học, thiếu tự tin, hoặc chỉ đơn giản là cần thêm thời gian.
\textit{(But slowness can come from many reasons: weak foundations, unsuitable methods, lack of confidence, or simply needing more time.)}

Nếu kết luận quá nhanh, em sẽ tự khóa mình lại.
\textit{(If you conclude too quickly, you lock yourself in.)}

Chậm không phải là dốt; rất nhiều người đi chậm nhưng đi bền.
\textit{(Slow is not stupid; many people move slowly but steadily.)}
\end{truyen}

\begin{baihoc}
Tốc độ không nói hết giá trị và khả năng của một người học.
\textit{(Speed does not fully define a learner's worth or ability.)}
\end{baihoc}

\begin{ghinhoanh}
\textbf{Short English to remember:}

Slow progress is still progress.

\end{ghinhoanh}

\begin{tuhocnhanh}
1. Vì sao chậm không đồng nghĩa với dốt? \textit{Why is being slow not the same as being stupid?}

2. Học chậm có thể đến từ những lý do nào? \textit{What can cause slower learning?}

3. Em có đang đánh giá mình quá nặng vì tốc độ không? \textit{Are you judging yourself too harshly because of pace?}
\end{tuhocnhanh}
\ngancach

\section{Mình Có Thể Tiến Bộ}
\begin{truyen}{Mình Có Thể Tiến Bộ}{You Can Improve}
\chuhoa{Đ}{iều rất quan trọng học sinh cần hiểu sớm là mình không đứng yên mãi như hiện tại.}
\textit{(One important thing students should understand early is that they are not frozen in their current state.)}

Con người có thể tiến bộ nhờ học, sửa, luyện và bền bỉ.
\textit{(People can improve through learning, correction, practice, and persistence.)}

Tiến bộ có thể chậm, nhưng điều đó khác hoàn toàn với việc không thể tiến bộ.
\textit{(Progress may be slow, but that is completely different from being impossible.)}

Tin vào khả năng tiến bộ giúp em có lý do để tiếp tục.
\textit{(Believing in improvement gives you reason to continue.)}
\end{truyen}

\begin{baihoc}
Không cần hoàn hảo ngay; điều cần là đừng từ chối khả năng lớn lên của mình.
\textit{(You do not need instant perfection; what matters is not denying your ability to grow.)}
\end{baihoc}

\begin{ghinhoanh}
\textbf{Short English to remember:}

Growth remains possible.

\end{ghinhoanh}

\begin{tuhocnhanh}
1. Vì sao học sinh cần tin mình có thể tiến bộ? \textit{Why do students need to believe they can improve?}

2. Tiến bộ thường đến qua những gì? \textit{What usually leads to growth?}

3. Em muốn tiến bộ nhất ở điểm nào? \textit{Where do you most want to improve?}
\end{tuhocnhanh}
\ngancach

\section{Sai Lầm Là Bình Thường}
\begin{truyen}{Sai Lầm Là Bình Thường}{Mistakes Are Normal}
\chuhoa{S}{ai là chuyện rất quen trong quá trình học và lớn lên.}
\textit{(Mistakes are a familiar part of learning and growing.)}

Nếu học sinh coi sai lầm như một điều quá ghê gớm, em sẽ dễ sợ học, sợ thử và sợ bắt đầu.
\textit{(If students treat mistakes as something terrible, they easily become afraid of learning, trying, and beginning.)}

Hiểu rằng sai là bình thường giúp em nhẹ hơn với bản thân.
\textit{(Understanding that mistakes are normal helps you be gentler with yourself.)}

Từ đó, em mới có thể sửa một cách đàng hoàng.
\textit{(From there, you can begin to correct things properly.)}
\end{truyen}

\begin{baihoc}
Sai lầm không làm em kém giá trị; nó chỉ nhắc rằng em còn đang học.
\textit{(Mistakes do not reduce your worth; they simply remind you that you are still learning.)}
\end{baihoc}

\begin{ghinhoanh}
\textbf{Short English to remember:}

Normal mistakes can still teach important lessons.

\end{ghinhoanh}

\begin{tuhocnhanh}
1. Vì sao sai lầm là bình thường? \textit{Why are mistakes normal?}

2. Coi sai quá nặng gây hại gì? \textit{What harm comes from treating mistakes too heavily?}

3. Em có đang quá khó với lỗi của mình không? \textit{Are you being too hard on your own mistakes?}
\end{tuhocnhanh}
\ngancach

\section{Tự Tin Có Thể Luyện}
\begin{truyen}{Tự Tin Có Thể Luyện}{Confidence Can Be Trained}
\chuhoa{N}{hiều học sinh nghĩ tự tin là thứ ai có sẵn thì có, ai không có thì chịu.}
\textit{(Many students think confidence is something you either naturally have or never will.)}

Thật ra, tự tin thường lớn dần từ những lần chuẩn bị tốt, những lần thử, và những lần vượt qua nỗi ngại nhỏ.
\textit{(In reality, confidence often grows through preparation, attempts, and small acts of courage.)}

Nó không phải phép màu xuất hiện một đêm.
\textit{(It is not a miracle that appears overnight.)}

Tự tin là một kỹ năng có thể tập được.
\textit{(Confidence is a skill that can be trained.)}
\end{truyen}

\begin{baihoc}
Đừng chờ tự tin rồi mới làm; nhiều khi chính việc làm tử tế mới tạo ra tự tin.
\textit{(Do not wait for confidence before acting; often good action is what builds confidence.)}
\end{baihoc}

\begin{ghinhoanh}
\textbf{Short English to remember:}

Confidence often grows through action.

\end{ghinhoanh}

\begin{tuhocnhanh}
1. Vì sao tự tin có thể luyện được? \textit{Why can confidence be trained?}

2. Nó thường lớn lên từ những gì? \textit{What does confidence usually grow from?}

3. Em có thể luyện tự tin ở việc nào trước? \textit{Where can you begin training confidence?}
\end{tuhocnhanh}
\ngancach

\section{Thay Đổi Là Có Thể}
\begin{truyen}{Thay Đổi Là Có Thể}{Change Is Possible}
\chuhoa{C}{ó học sinh rất sớm đã nghĩ: tính em thế rồi, chắc không đổi được đâu.}
\textit{(Some students think very early: this is just who I am, I probably cannot change.)}

Nhưng con người thay đổi rất nhiều qua thói quen, cách nghĩ và những lựa chọn lặp đi lặp lại.
\textit{(But people change greatly through habits, thought patterns, and repeated choices.)}

Không phải điều gì cũng đổi nhanh, nhưng rất nhiều thứ vẫn có thể tốt hơn.
\textit{(Not everything changes quickly, but many things can still improve.)}

Hiểu điều này giúp em bớt đóng khung chính mình.
\textit{(Understanding this keeps you from trapping yourself in a fixed box.)}
\end{truyen}

\begin{baihoc}
Chậm thay đổi không có nghĩa là không thể thay đổi.
\textit{(Slow change does not mean impossible change.)}
\end{baihoc}

\begin{ghinhoanh}
\textbf{Short English to remember:}

Change may be slow, but it is still real.

\end{ghinhoanh}

\begin{tuhocnhanh}
1. Vì sao nhiều học sinh tự đóng khung mình? \textit{Why do many students place themselves in fixed boxes?}

2. Con người thay đổi qua những gì? \textit{Through what do people change?}

3. Em muốn thay đổi điều gì trước? \textit{What do you want to change first?}
\end{tuhocnhanh}
\ngancach

\section{Tin Vào Bản Thân}
\begin{truyen}{Tin Vào Bản Thân}{Believe in Yourself}
\chuhoa{T}{in vào bản thân không có nghĩa là nghĩ mình hơn người khác.}
\textit{(Believing in yourself does not mean thinking you are above others.)}

Nó chỉ có nghĩa là em không vội phủ nhận giá trị và khả năng lớn lên của mình.
\textit{(It simply means you do not rush to deny your value and your ability to grow.)}

Một người không tin vào mình rất dễ bỏ cuộc trước cả khi thật sự thử đủ.
\textit{(A person who does not believe in themselves gives up too easily before really trying enough.)}

Tin vào bản thân là một điểm tựa rất cần cho tuổi học trò.
\textit{(Self-belief is an important support for student life.)}
\end{truyen}

\begin{baihoc}
Niềm tin đúng vào bản thân không làm em kiêu ngạo; nó giúp em không tự xóa mình đi.
\textit{(Healthy self-belief does not make you arrogant; it keeps you from erasing yourself.)}
\end{baihoc}

\begin{ghinhoanh}
\textbf{Short English to remember:}

Healthy self-belief supports steady growth.

\end{ghinhoanh}

\begin{tuhocnhanh}
1. Tin vào bản thân khác gì với kiêu ngạo? \textit{How is self-belief different from arrogance?}

2. Thiếu niềm tin vào mình dễ dẫn đến điều gì? \textit{What does a lack of self-belief often cause?}

3. Em cần nói với mình câu nào công bằng hơn? \textit{What fairer sentence do you need to tell yourself?}
\end{tuhocnhanh}
\ngancach